\section{Mathematical Model}
\label{appendix:dynamics}
This section derives the full nonlinear dynamics of the ball-on-arc system. The ball rolls on a convex arc with a central concave dip. The dip makes the equilibrium locally stable and lightly damped, a shallow restoring well, rather than open-loop unstable as in the classical ball-and-beam system. The well is narrow, however: with the parameters of Table~\ref{tab:sys_params}, gravity restores the ball only within about $0.0115$\,rad of center and beyond that the convex arc dominates and rolls the ball away. We use the Euler-Lagrange method under a pure-rolling assumption, accounting for both translational and rotational motion of the ball. The base ball-and-arc formulation follows~\cite{ho2010smc}; we extend it with the central concave dip (the $d\,e^{-k\theta^2}$ term below), rolling and viscous friction, and a velocity-input actuator model (Section~\ref{sec:velocity_input}).

\subsection{Coordinates and Kinematics}

We define the generalized coordinates as $x$ (cart position) and $\theta$ (ball angle measured from the arc's vertical centerline, positive to the right), consistent with the platform schematic in the main paper. The ball position relative to the cart is
\begin{equation}
X_{\mathrm{rel}} = R \sin\theta, \quad
Y_{\mathrm{rel}} = R \cos\theta - d e^{-k\theta^2},
\end{equation}
where $R$ is the arc radius, $d$ is the dip depth, and $k = R^2/(2\sigma^2)$ defines the Gaussian dip width. The term $d e^{-k\theta^2}$ represents the concave dip that lowers the surface by depth $d$ at the arc center ($\theta = 0$) and smoothly transitions to the nominal convex arc profile as $|\theta|$ increases. In inertial coordinates,
\begin{equation}
X = x + X_{\mathrm{rel}}, \quad Y = Y_{\mathrm{rel}}.
\end{equation}

The corresponding velocity components are
\begin{equation}
\label{eq:XYcoordVel}
\dot X = \dot x + R \cos\theta \, \dot\theta, \quad
\dot Y = (-R \sin\theta + a(\theta)) \dot\theta,
\end{equation}
with $a(\theta) = 2 d k \theta e^{-k\theta^2}$ capturing the velocity contribution from the dip geometry.

Assuming pure rolling without slip, the ball's angular velocity is
\begin{equation}
\omega = \frac{R}{r} \dot\theta,
\end{equation}
where $r$ is the ball radius. This holonomic constraint $r \varphi = R \theta$ relates the ball's rotational angle $\varphi$ to its position along the arc.

\subsection{Energy Formulation}

\subsubsection{Translational Kinetic Energy}
The cart moves along the horizontal track:
\begin{equation}
T_\text{cart} = \frac{1}{2} M \dot x^2.
\end{equation}

The ball contributes translational kinetic energy in both the horizontal and vertical directions:
\begin{equation}
T_\text{ball,trans} = \frac{1}{2} m (\dot X^2 + \dot Y^2).
\end{equation}

Using \eqref{eq:XYcoordVel}, we expand:
\begin{equation}
\begin{split}
\dot X^2 + \dot Y^2 &= \dot x^2 + 2 \dot x R \cos\theta \, \dot\theta \\[2pt]
&\quad + \bigl[ (R\cos\theta)^2 + (-R\sin\theta + a(\theta))^2 \bigr] \dot\theta^2.
\end{split}
\end{equation}

Combining cart and ball translational contributions:
\begin{equation}
T_\text{trans} = \frac{1}{2}(M+m)\dot x^2 + m R \cos\theta \, \dot x \dot\theta + \frac{1}{2} m C(\theta) \dot\theta^2,
\end{equation}
with
\begin{equation}
C(\theta) = (R\cos\theta)^2 + (-R\sin\theta + a(\theta))^2.
\end{equation}

\subsubsection{Rotational Kinetic Energy}
The ball rolls without slipping along the arc:
\begin{equation}
T_\text{rot} = \frac{1}{2} I \omega^2 = \frac{1}{2} I \left(\frac{R}{r}\right)^2 \dot\theta^2,
\end{equation}
where $I$ is the moment of inertia of the ball and $r$ its radius.

\subsubsection{Effective Ball Contribution}
Combining translational and rotational contributions along the arc:
\begin{equation}
C_\mathrm{eff}(\theta) = C(\theta) + \frac{I}{m} \frac{R^2}{r^2}, \quad
T_\text{ball} = \frac{1}{2} m C_\mathrm{eff}(\theta) \dot\theta^2.
\end{equation}

\subsubsection{Total Kinetic and Potential Energy}
\begin{equation}
\boxed{
T = \frac{1}{2}(M+m)\dot x^2 + m R \cos\theta \, \dot x \dot\theta + \frac{1}{2} m C_\mathrm{eff}(\theta) \dot\theta^2
}
\end{equation}
\begin{equation}
\boxed{V(\theta) = mg \big(R \cos\theta - d e^{-k\theta^2} \big).}
\end{equation}

\subsection{Equations of Motion}
The Euler--Lagrange equation for a generalized coordinate $q_i$ with generalized force $Q_i$ is:
\begin{equation}
\frac{d}{dt} \left( \frac{\partial L}{\partial \dot{q}_i} \right) - \frac{\partial L}{\partial q_i} = Q_i,
\label{eq:euler-lagrange}
\end{equation}
where $L = T - V$ is the Lagrangian.

The derivatives used in the formulation are defined as
\begin{equation}
C'(\theta) = \frac{d C_\mathrm{eff}}{d\theta}, \quad
a'(\theta) = \frac{d a(\theta)}{d\theta} = 2 d k e^{-k\theta^2} (1 - 2 k \theta^2).
\end{equation}

\paragraph{Cart coordinate $x$}  
The horizontal force applied is $Q_x = u$. Partial derivatives of the Lagrangian:
\begin{equation}
\begin{split}
\frac{\partial L}{\partial \dot{x}} &= (M+m)\dot{x} + m R \cos\theta \, \dot\theta, \\
\frac{\partial L}{\partial x} &= 0.
\end{split}
\end{equation}

Time derivative:
\begin{equation}
\begin{split}
\frac{d}{dt} \left( \frac{\partial L}{\partial \dot{x}} \right) 
&= (M+m)\ddot{x} + m R \left(-\sin\theta \, \dot\theta^2 + \cos\theta \, \ddot\theta \right).
\end{split}
\end{equation}


The Euler--Lagrange equation gives the following:
\begin{equation}
\label{eq:x_final_eq}
\boxed{(M+m)\ddot{x} - m R \sin\theta \, \dot\theta^2 + m R \cos\theta \, \ddot\theta = u.}
\end{equation}

\paragraph{Ball coordinate $\theta$}  
The generalized force is zero, $Q_\theta = 0$. Partial derivatives of the Lagrangian are
\begin{equation}
\begin{split}
\frac{\partial L}{\partial \dot{\theta}} &= m R \cos\theta \, \dot{x} + m C_\mathrm{eff}(\theta) \, \dot\theta, \\
\frac{\partial L}{\partial \theta} &= - m R \sin\theta \, \dot{x} \dot{\theta}
  + \tfrac{1}{2} m C_\mathrm{eff}'(\theta) \, \dot\theta^2
  - m g \bigl(a(\theta) - R\sin\theta\bigr).
\end{split}
\end{equation}

Time derivative:
\begin{equation}
\begin{split}
\frac{d}{dt} \left( \frac{\partial L}{\partial \dot{\theta}} \right) 
&= m R (-\sin\theta \, \dot\theta \, \dot{x} + \cos\theta \, \ddot{x}) \\[2pt]
&\quad + m \left( C'(\theta) \dot\theta^2 + C_\mathrm{eff}(\theta) \ddot{\theta} \right).
\end{split}
\end{equation}

The Euler--Lagrange equation simplifies to:
\begin{equation}
\label{eq:theta_final_eq}
\boxed{
R \cos\theta\,\ddot{x}
+ C_\mathrm{eff}(\theta)\,\ddot{\theta}
+ \tfrac{1}{2}C'(\theta)\dot{\theta}^2
+ g\bigl(a(\theta) - R\sin\theta\bigr) = 0
}
\end{equation}

\subsection{Matrix Form}
Coupled second-order dynamics from Eqs.~\eqref{eq:x_final_eq} and \eqref{eq:theta_final_eq} can be expressed compactly in matrix form:
\begin{equation}
\label{eq:matrix_eom}
\begin{split}
&\begin{bmatrix} 
M+m & m R \cos\theta \\[1ex] 
R \cos\theta & C_{\mathrm{eff}}(\theta)
\end{bmatrix}
\begin{bmatrix} \ddot x \\ \ddot\theta \end{bmatrix} = \\[4pt]
&\qquad
\begin{bmatrix} 
u + m R \sin\theta \,\dot\theta^2 \\[1ex] 
- \tfrac{1}{2} C_{\mathrm{eff}}'(\theta)\,\dot\theta^2 
- g\big(-R \sin\theta + a(\theta)\big)
\end{bmatrix}.
\end{split}
\end{equation}

This matrix form makes it easier to see the coupling between the cart and ball. For controller design, we convert the system into a first-order state-space form. 

\subsection{First-Order State-Space Form}

To facilitate controller design and linearization, we convert the second-order dynamics in Eq.~\eqref{eq:matrix_eom} to a first-order form. We define the state vector as
\begin{equation}
\label{eq:state_vector}
x = 
\begin{bmatrix}
x_\text{cart} \\ \dot{x}_\text{cart} \\ \theta \\ \dot{\theta}
\end{bmatrix} \in \mathbb{R}^4, \quad
u \in \mathbb{R}.
\end{equation}

The dynamics can then be written compactly as
\begin{equation}
\label{eq:first_order_dynamics}
\dot{x} = f(x,u) =
\begin{bmatrix}
\dot{x}_\text{cart} \\
\ddot{x}_\text{cart} (x, \dot{x}, \theta, \dot{\theta}, u) \\
\dot{\theta} \\
\ddot{\theta} (x, \dot{x}, \theta, \dot{\theta})
\end{bmatrix},
\end{equation}
where $\ddot{x}_\text{cart}$ and $\ddot{\theta}$ are computed from Eq.~\eqref{eq:matrix_eom}.

\paragraph{Linearization}  
We linearize the system around an operating point $(x_0,u_0)$ using the Jacobians of $f$ with respect to the state and input:
\begin{equation}
\label{eq:jacobians}
A = \left. \frac{\partial f}{\partial x} \right|_{x_0, u_0}, \quad
B = \left. \frac{\partial f}{\partial u} \right|_{x_0, u_0}.
\end{equation}
These matrices form the basis for designing linear controllers.

\paragraph{Discrete-Time Approximation}  
For controllers that require discrete-time dynamics, the first-order system can be integrated over a timestep $\Delta t$ using zero-order hold (ZOH):
\begin{equation}
\label{eq:discrete_map}
x_{k+1} = F(x_k, u_k) \approx x_k + \int_{t_k}^{t_k + \Delta t} f(x(t), u_k) \, dt,
\end{equation}
where $F$ defines the discrete-time map. 

\subsection{Velocity-Input Plant Model}
\label{sec:velocity_input}

The physical platform uses a Beckhoff servo drive controlled via
the \texttt{MC\_MoveVelocity} function block, where the controller
commands a desired cart velocity~$v_\text{cmd}$ rather than a force.
Internally, the drive executes a 7-phase S-curve motion profile;
for modeling tractability, we approximate this as a first-order lag:
\begin{equation}
\ddot{x} = \frac{v_\text{cmd} - \dot{x}}{\tau_v},
\qquad \tau_v = 0.15~\text{s},
\label{eq:velocity_lag}
\end{equation}
where $\tau_v$ is identified empirically from open-loop
trajectory matching between simulation and hardware
(Appendix~\ref{appendix:simulation_validation}).
Because $\ddot{x}$ is now a known function of state and input,
the ball equation reduces to
\begin{equation}
\begin{split}
\ddot{\theta} &= \frac{1}{C_{\mathrm{eff}}(\theta)}\Bigl[
  -R\cos\theta\,\ddot{x}
  - \tfrac{1}{2}C_{\mathrm{eff}}'(\theta)\dot{\theta}^2 \\
  &\quad + g(R\sin\theta - a(\theta))
  - Q_\mathrm{rr} - Q_\mathrm{visc}
\Bigr],
\end{split}
\label{eq:ball_reduced}
\end{equation}
where the generalized friction terms are
\begin{align}
Q_\mathrm{rr}   &= \mu_\mathrm{rr}\, g\, R\cos\theta\,
                    \operatorname{sign}(\dot\theta),
                    \label{eq:Q_rr} \\
Q_\mathrm{visc} &= \beta_v\, \dot\theta.
                    \label{eq:Q_visc}
\end{align}
with rolling-friction coefficient
$\mu_\mathrm{rr} = 0.025$ and viscous damping
$\beta_v = 0.009~\text{m}^2/\text{s}$. These two are nominal estimates
rather than precisely identified values; domain randomization samples
them over $[0.010, 0.040]$ and $[0.003, 0.015]$ respectively during
training. In simulation the $\operatorname{sign}(\dot\theta)$ in
Eq.~\eqref{eq:Q_rr} is replaced by a smoothed
$\tanh(\dot\theta/\varepsilon)$ with $\varepsilon = 0.005$\,rad/s to
avoid chattering at $\dot\theta = 0$.
The physical parameters are taken from \texttt{balancer/core/params.py};
the actuation limit $v_\mathrm{max}$ and the 20\,Hz control rate come from
\texttt{balancer/hardware/constants.py}. All are summarized in
Table~\ref{tab:sys_params}.

\begin{table}[!tb]
\centering
\small
\setlength{\tabcolsep}{4pt}
\caption{Nominal system parameters.}
\label{tab:sys_params}
\begin{tabular}{llp{2.9cm}}
\toprule
\textbf{Parameter} & \textbf{Value} & \textbf{Description} \\
\midrule
$g$                & 9.81\,m/s$^2$  & Gravitational acceleration \\
$R$                & 2.101\,m        & Arc radius \\
$m$              & 0.024\,kg       & Ball mass (24\,g) \\
$r$              & 0.009\,m        & Ball radius (9\,mm) \\
$M$              & 0.351\,kg       & Cart mass \\
$\tau_v$           & 0.15\,s         & Velocity-lag time constant \\
$\mu_\mathrm{rr}$  & 0.025           & Ball rolling friction (est.) \\
$\beta_v$          & 0.009\,m$^2$/s  & Ball viscous damping (est.) \\
$d$                & 0.002\,m        & Dip depth (2\,mm) \\
$\mathrm{FWHM}$    & 0.020\,m        & Dip width ($\sigma=$\,FWHM$/2.355$) \\
$v_\mathrm{max}$   & 0.9\,m/s        & Cart velocity-command limit \\
$T_s$              & 0.05\,s         & Control timestep (20\,Hz) \\
\bottomrule
\end{tabular}
\end{table}

All reported controllers and the benchmark environments use this
velocity-input plant; the drive's underlying jerk-limited S-curve
profile is also implemented and drives the development-regime results
(Appendix~\ref{appendix:rl_fairness}), while the force-input formulation
(Eq.~\eqref{eq:matrix_eom}) is retained for derivation
reference only.  Both the linearized velocity-input system
(used by LQR/MPC) and the nonlinear RK4 integration
(used by NMPC) incorporate $\tau_v$ explicitly.

\subsection{Remarks}
We treat the ball as a solid spherical bead, so its moment of inertia is
$I = \tfrac{2}{5} m r^2$. LQR and MPC use this model linearized around the arc
center, while NMPC and the training simulator integrate the full nonlinear
form; in either case the model carries the coupling between cart and ball
along with the effect of the central concave dip.


% ============================================================
%  SECTION IV - EXPERIMENTS
% ============================================================
